\chapter{Data modeling}

Data is the backend's most important interface.
Schema decisions are long-lived and difficult to change.

\section{Entities and relationships}
Start with the core entities and how they relate.
Use a small number of tables with clear ownership.

\section{Constraints}
Use database constraints for uniqueness and referential integrity.
Do not rely solely on application checks.

\section{Constraint examples}
Constraints encode invariants directly in the database.
\begin{lstlisting}[language=SQL]
ALTER TABLE tasks
  ADD CONSTRAINT tasks_status_check
  CHECK (status IN ('open', 'in_progress', 'done'));
\end{lstlisting}

\begin{lstlisting}[language=SQL]
CREATE TABLE tasks (
  id TEXT PRIMARY KEY,
  project_id TEXT NOT NULL REFERENCES projects(id),
  title TEXT NOT NULL,
  status TEXT NOT NULL,
  created_at TIMESTAMP NOT NULL
);
\end{lstlisting}

\section{Indexes}
Indexes are a performance contract.
Add indexes for your query patterns, not for every column.

\section{Normalization and denormalization}
Normalize to reduce duplication and ensure correctness.
Denormalize when read performance is more important than write complexity.

\section{Data ownership}
Each table should have a clear owner in the domain model.
Ownership determines which service or module is allowed to write the data.

\section{Query patterns}
Design the schema around the queries you will actually run.
If you regularly filter tasks by project and status, index both fields.

\section{Example query patterns}
\begin{tabularx}{\textwidth}{@{}lX@{}}
\toprule
Query & Suggested index \\
\midrule
List tasks by project and status & (\code{project\_id}, \code{status}, \code{created\_at}) \\
Fetch task by id & (id) \\
List recent activity by task & (\code{task\_id}, \code{created\_at}) \\
\bottomrule
\end{tabularx}

\section{Migrations}
Schema changes should be backward compatible.
Add columns before you use them, and remove them later.

\section{Migration example}
\begin{lstlisting}[language=SQL]
BEGIN;
ALTER TABLE tasks ADD COLUMN priority INTEGER DEFAULT 0;
CREATE INDEX idx_tasks_project_status ON tasks(project_id, status);
COMMIT;
\end{lstlisting}

\section{Soft deletes}
Soft deletes keep history but complicate queries.
If you need history, consider an audit table instead of a \code{deleted} flag everywhere.

\section{Audit trails}
Write audit logs for critical actions like permissions changes and deletions.
These logs are essential for debugging and compliance.

\section{Data retention}
Not all data should live forever.
Define retention policies for logs, activity, and temporary data.

\section{Backfills}
If you add a new column, you may need to backfill old data.
Plan backfills as separate jobs so they do not block user-facing requests.

\section{Transactions and isolation}
Use transactions to keep multi-step changes consistent.
Keep them short and avoid long-running queries inside a transaction.

\section{Read replicas}
Read replicas improve scale but introduce replication lag.
If you require read-after-write behavior, read from the primary.

\section{Explain plans}
Use query explain plans to understand performance.
Add indexes only when they improve the actual query plan.

\section{Integrity checks}
Schedule periodic integrity checks to catch data drift.
These are especially important after large migrations or backfills.

\section{Write amplification}
Denormalization and heavy indexing increase write cost.
Measure write latency before adding more indexes.

\section{Schema evolution}
Use additive changes first, then migrate clients, then remove old fields.
This minimizes downtime and unexpected client failures.

\section{Partitioning}
Partitioning can keep large tables manageable.
Choose a partition key that matches query patterns, such as project id or date.

\section{Backups and recovery}
Backups are part of the data model.
If you cannot restore data reliably, your model is not production-ready.

\section{Data locality}
Place data close to the services that use it most.
Cross-region reads add latency and can create consistency surprises.

\section{Case study: activity log growth}
Activity logs grow faster than core entities.
If you store them in the same table as tasks, reads and writes will slow down.
A dedicated activity table with time-based partitioning keeps growth manageable.

\section{Data validation at rest}
Use constraints and types to prevent invalid data from being stored.
If your database supports it, use check constraints for enums and ranges.

\section{Null handling}
Nulls are easy to add and hard to reason about.
Define when null is allowed and what it means.
If a field is required for business logic, make it \code{NOT NULL}.

\section{Data lifecycle and classification}
Not all data is equally sensitive.
Classify fields early:
\begin{itemize}[nosep]
  \item public metadata (project name)
  \item internal operational data (status transitions)
  \item sensitive data (user identifiers, billing details)
\end{itemize}

Classification drives encryption, access controls, and retention.

\section{Hot and cold data}
Separate hot data from cold data to keep critical queries fast.
For example, keep active tasks in a primary table and archive completed tasks.
This keeps indexes small and improves cache hit rates.

\section{Reference data}
Some data changes rarely, such as status codes or categories.
Model it as reference tables and enforce foreign keys.
This reduces duplication and keeps validation consistent.

\section{Sharding criteria}
If the system grows, sharding may be necessary.
Pick a key that keeps related data together, such as project id.
Avoid sharding on fields that cause cross-shard queries for core workflows.

\section{Event tables}
Event tables store a log of changes instead of only the current state.
They are useful for audits and analytics.
If you use event tables, plan retention and compaction to control size.

\section{Change data capture}
Some systems need real-time streams of data changes.
Use change data capture or an outbox to publish changes to other services.
This keeps integration logic out of the primary request path.

\section{Precision and types}
Choose data types that match the domain:
\begin{itemize}[nosep]
  \item use \code{DECIMAL} for money
  \item use \code{TIMESTAMP WITH TIME ZONE} for time
  \item avoid free-text for enums
\end{itemize}

\section{Testing with realistic data}
Synthetic data that is too small hides problems.
Keep a staging dataset that reflects production volume and skew.
This helps validate query plans and migration timing.

\section{Hotspot avoidance}
Some keys receive a disproportionate amount of traffic.
If a single project or user dominates traffic, you can create hotspots.
Mitigate hotspots with sharding, caching, or dedicated queues.

\section{Index maintenance}
Indexes speed reads but slow writes.
Periodically review unused or low-value indexes.
Drop indexes that no longer align with real query patterns.

\section{Data fixes and runbooks}
Production data will need fixes.
Write a runbook for safe data changes:
\begin{itemize}[nosep]
  \item define the scope and affected records
  \item run changes in small batches
  \item verify metrics and constraints after the change
\end{itemize}

\section{PII and compliance}
Personally identifiable information requires special handling.
Apply stricter access controls and limit exposure in logs.
If you must delete user data, plan for hard deletes or anonymization.

\section{Analytics schema}
Operational schemas are not always ideal for analytics.
Use a separate analytics schema or warehouse for reporting.
This protects the primary database from heavy reporting queries.

\section{Time-series data}
Time-series data like metrics and events grows quickly.
Use time-based partitioning and retention policies.
Avoid mixing long-term time-series data with transactional tables.

\section{Foreign keys and deletes}
Define delete behavior explicitly.
Use \code{ON DELETE RESTRICT} for critical data and \code{ON DELETE CASCADE} only when safe.
Unexpected cascades can delete far more data than intended.

\section{Version columns}
Version columns make data changes traceable.
They support optimistic locking and conflict detection.
If you update a record, increment its version and record the timestamp.

\section{Multi-tenant schemas}
For multi-tenant data, choose a strategy early:
\begin{itemize}[nosep]
  \item shared tables with a \code{tenant\_id} column
  \item separate schemas per tenant
  \item separate databases for large tenants
\end{itemize}

Each strategy changes operational complexity and isolation.

\section{Data quality monitoring}
Data quality is an operational metric.
Track null rates, duplicate keys, and constraint violations.
If you do not measure data quality, it will regress silently.

\section{Archival strategy}
Define what happens when data ages out.
Archive to cheaper storage or delete after the retention window.
Make the archival process visible and reversible.

\section{Access patterns and tooling}
The way you access data shapes the schema.
If you use an ORM, understand its query patterns and add indexes accordingly.
If you write raw SQL, standardize queries and avoid ad-hoc joins.

\section{Partitioning example}
Partitioning is only useful when it matches access patterns.
For a task system, time-based partitioning is common for activity logs:

\begin{tabularx}{\textwidth}{@{}lX@{}}
\toprule
Table & Partition key \\
\midrule
activity\_log & month by created\_at \\
tasks & project\_id or created\_at \\
\bottomrule
\end{tabularx}

\section{RPO and RTO}
Backups should meet business objectives.
Define recovery point objective (RPO) and recovery time objective (RTO) per dataset.
If you cannot meet them, adjust either the tooling or the expectations.

\section{Materialized views}
If you need expensive aggregates, consider materialized views.
They trade storage and refresh complexity for faster reads.
Document refresh schedules and acceptable staleness.

\section{Data governance}
Define who can change schema, who can access production data, and how audits are performed.
Governance should be lightweight but explicit to avoid accidental exposure.

\section{Schema documentation}
Schema docs should live near the code.
At minimum, document table purpose, key fields, and ownership.
Good documentation speeds onboarding and reduces misuses.

\section{Backfill plan example}
Backfills should be treated like data migrations:
\begin{itemize}[nosep]
  \item add the new column with a default
  \item deploy code that reads both old and new fields
  \item run a backfill in batches with progress tracking
  \item validate counts and constraints
  \item remove the old field after verification
\end{itemize}

\section{Migration phases}
Large migrations work best in explicit phases:

\begin{tabularx}{\textwidth}{@{}lX@{}}
\toprule
Phase & Goal \\
\midrule
Prepare & add schema changes and deploy dual-read code \\
Backfill & migrate existing records in controlled batches \\
Cutover & switch writers to the new schema \\
Cleanup & remove deprecated fields and indexes \\
\bottomrule
\end{tabularx}

\section{Validation queries}
After a migration, verify the data with simple queries:
\begin{lstlisting}[language=SQL]
SELECT COUNT(*) FROM tasks WHERE priority IS NULL;
SELECT COUNT(*) FROM tasks WHERE project_id IS NULL;
SELECT status, COUNT(*) FROM tasks GROUP BY status;
\end{lstlisting}

\section{Sampling for checks}
Full scans can be expensive.
Use sampling for daily checks and full scans for scheduled audits.
Make the sampling strategy part of the runbook.

\section{Indexing case study}
Assume most reads are \code{GET /tasks?project=id}.
An index on \code{project\_id} is required, but adding \code{status} can improve filtering.
Measure the query plan before and after to confirm the index is used.

\section{Retention and purge jobs}
Retention is not just a policy, it is a job.
Schedule purge jobs and track their runtime and impact.
If a purge job takes hours, it will compete with production traffic.

\section{Data drift monitoring}
Data can drift after migrations and backfills.
Monitor for drift with periodic checks:
\begin{itemize}[nosep]
  \item compare counts between old and new columns
  \item validate enums against reference tables
  \item sample rows for manual verification
\end{itemize}

\section{Schema example}
Keeping schema examples in the book makes the discussion concrete:
\begin{lstlisting}[language=SQL]
CREATE TABLE projects (
  id TEXT PRIMARY KEY,
  name TEXT NOT NULL,
  owner_id TEXT NOT NULL,
  created_at TIMESTAMP NOT NULL
);

CREATE TABLE tasks (
  id TEXT PRIMARY KEY,
  project_id TEXT NOT NULL REFERENCES projects(id),
  title TEXT NOT NULL,
  status TEXT NOT NULL,
  created_at TIMESTAMP NOT NULL
);
\end{lstlisting}

\section{Index and query example}
Indexes should map directly to hot queries:
\begin{lstlisting}[language=SQL]
CREATE INDEX idx_tasks_project_status
ON tasks(project_id, status, created_at);

SELECT id, title
FROM tasks
WHERE project_id = 'proj_7' AND status = 'open'
ORDER BY created_at DESC
LIMIT 50;
\end{lstlisting}

\section{Audit table example}
Audit tables preserve history without bloating primary tables:
\begin{lstlisting}[language=SQL]
CREATE TABLE task_audit (
  id TEXT PRIMARY KEY,
  task_id TEXT NOT NULL REFERENCES tasks(id),
  action TEXT NOT NULL,
  actor_id TEXT NOT NULL,
  created_at TIMESTAMP NOT NULL
);
\end{lstlisting}

\section{Data dictionary example}
A short data dictionary clarifies meaning and ownership:

\begin{tabularx}{\textwidth}{@{}lX@{}}
\toprule
Field & Meaning \\
\midrule
projects.owner\_id & user who owns the project \\
tasks.status & lifecycle state of a task \\
tasks.created\_at & time the task was created \\
task\_audit.action & action performed on the task \\
task\_audit.actor\_id & user who performed the action \\
\bottomrule
\end{tabularx}

\begin{patternbox}{Pattern: model change as an interface change}
If a schema change breaks clients, treat it like a breaking API change.
\end{patternbox}

\begin{checklistbox}{Checklist: data modeling}
\begin{itemize}[nosep]
  \item Use constraints for integrity.
  \item Index the queries you actually run.
  \item Migrate in small, reversible steps.
\end{itemize}
\end{checklistbox}
